\documentclass[11pt,a4paper]{article}
\usepackage{fontspec}
\usepackage[spanish,es-noshorthands]{babel}
\usepackage[margin=2.3cm]{geometry}
\usepackage{booktabs,array,tabularx}
\usepackage{enumitem}
\usepackage{amsmath}
\usepackage{xcolor}
\usepackage[hidelinks]{hyperref}
\usepackage{titlesec}
\usepackage{parskip}
\usepackage{needspace}

\definecolor{acento}{RGB}{30,70,120}
\titleformat{\section}{\large\bfseries\color{acento}}{\thesection.}{0.5em}{}
\titlespacing*{\section}{0pt}{14pt}{4pt}
\setlist{itemsep=1pt,topsep=2pt}
\newcolumntype{L}{>{\raggedright\arraybackslash}X}

% Caja corta al inicio de cada punto: comentario y respuesta en una línea
\newcommand{\encorto}[2]{\par\noindent\colorbox{acento!7}{\parbox{\dimexpr\linewidth-2\fboxsep}{%
\small\textbf{Comentario:} #1\\[1pt]\textbf{Respuesta corta:} #2}}\par\smallskip}
\newcommand{\paso}[1]{\par\smallskip\needspace{5\baselineskip}\noindent\textbf{\color{acento}#1}\par}

\title{\vspace{-1.5cm}\textbf{Respuesta a los comentarios sobre el relatorio del 16/10}\\[2pt]
\large SAR biestático complementario para sensado subsuperficial}
\author{David Atencia}
\date{20 de octubre de 2026}

\begin{document}
\maketitle
\vspace{-0.6cm}

Desde la versión que le entregué, el relatorio se actualizó (versión del 18/10, 57~pp.).
Ahora incluye el Exp.~C2 y corrige los pasajes que presentaban TE a 90° como el mejor
resultado del estudio. Abajo respondo cada comentario y, al final, dejo tres preguntas para
la próxima reunión.

\section*{Resumen}

\begin{tabularx}{\linewidth}{@{}>{\raggedright\arraybackslash}p{3.3cm}LL@{}}
\toprule
Punto & Respuesta & Próximo paso \\
\midrule
1. Validar la trayectoria &
De acuerdo: se eligió sólo con el blanco puntual. &
Reevaluar los 10--15 mejores candidatos con objetos reales y clutter. \\
\addlinespace
2. Nulo TE a 90° &
Era un artefacto del blanco puntual. El nulo útil es el VV en TM a 135°--140°. &
Verificarlo con un modelo de superficie de orden superior. \\
\addlinespace
3. Combinar canales &
De acuerdo: pesos iguales desperdician el canal bi. &
Comparar cuatro métodos de combinación. \\
\addlinespace
4. Calibración &
Hipótesis de trabajo: el desfase entre canales es constante en la colecta. &
Evaluación previa de la deriva esperada; simular sólo lo que supere la tolerancia. \\
\addlinespace
5. Aperturas parciales y suelos húmedos &
Faltan aperturas parciales. La pérdida en suelo húmedo no depende de la geometría. &
Repetir con arcos parciales y calcular la profundidad máxima por suelo. \\
\bottomrule
\end{tabularx}

\medskip
\noindent\textbf{Términos.}
\emph{Mono}: imagen del transceptor solo. \emph{Bi}: imagen del receptor separado.
\emph{Fusión}: suma de las dos imágenes.
$\Delta\phi$: separación angular entre Tx y Rx alrededor de la escena.
SCR: relación entre el eco del blanco y el clutter del suelo (más alto es mejor).
ISLR: energía de los lóbulos laterales de la imagen (más bajo es mejor).
SO(3): orientaciones aleatorias del objeto en 3D.

% ════════════════════════════════════════════════════════════════════
\section{Validación de la trayectoria propuesta}

\encorto{la trayectoria se eligió con un blanco puntual y no se validó con objetos reales.}
{tiene razón; hay que reevaluarla.}

La trayectoria del Exp.~F se eligió por la calidad de imagen de un punto (ISLR). Las
pruebas con objetos reales en orientaciones aleatorias (Exp.~E) se hicieron con
separaciones fijas (45°, 90°, 135° y 180°), no con el barrido de 60° a 150° de esa
trayectoria.

\paso{Qué propongo}
Reevaluar los 10 a 15 mejores candidatos, no sólo el elegido, con la misma prueba para
todos: objetos reales con orientación aleatoria, Rx de doble polarización y suelo rugoso
con clutter. El criterio sería el contraste típico (mediana) y la probabilidad de que la
fusión mejore la imagen mono. La elección puede cambiar: con objetos reales, la zona de
135° a 140° en TM vale más de lo que mostraba el punto (ver punto~2).

La prueba ya está implementada (Exp.~E2); falta aplicarla a las trayectorias del Exp.~F.

% ════════════════════════════════════════════════════════════════════
\section{El nulo TE a 90°}

\encorto{el nulo de clutter en TE a 90° puede ser un artefacto del modelo.}
{lo es, aunque por otra razón: el blanco puntual no tiene polarización. El nulo que sí
sirve es otro, en TM a 135°--140°.}

Después de entregar el relatorio ejecuté el Exp.~C2, que repite la prueba con una esfera y
un cubo metálicos. A 90° en TE el suelo deja de devolver clutter, pero el objeto real
tampoco devuelve eco. Sólo el punto escalar, que no tiene polarización, mostraba la mejora.

\begin{center}
\begin{tabular}{@{}lcc@{}}
\toprule
SCR en TE [dB] & Mono & Bi 90° \\
\midrule
Punto escalar & 3.9 & \textbf{21.7} \\
Esfera & 2.3 & $-4.8$ \\
Cubo, orientación aleatoria & $-8.8$ & $-9.0$ \\
\bottomrule
\end{tabular}
\end{center}

Lo que sí funciona con objetos reales es el nulo en TM con $\Delta\phi$ entre 135° y 140°
(un ``Brewster biestático'' del clutter). Ahí la imagen bi sola mejora el SCR frente a la
mono en 8 a 11~dB para la esfera y en 18 a 20~dB para el cubo con orientación aleatoria.

\begin{center}
\begin{tabular}{@{}lcccc@{}}
\toprule
SCR en TM [dB] & Mono & Bi 135° & Bi 140° & Fusión 140° \\
\midrule
Esfera & $-0.9$ & 7.5 & \textbf{9.8} & $-0.9$ \\
Cubo nominal & $-9.5$ & 14.7 & \textbf{20.1} & 4.7 \\
Cubo, orientación aleatoria & $-12.0$ & 6.2 & \textbf{8.3} & $-4.2$ \\
\bottomrule
\end{tabular}
\end{center}

\paso{Qué propongo}
Este nulo también sale del modelo de superficie más simple (SPM de primer orden), así que
su duda sigue en pie. Lo verificaría con:
\begin{itemize}
  \item un modelo de superficie de orden superior (IEM/AIEM o SPM de segundo orden);
  \item polarización cruzada de la superficie y del blanco, que hoy no se modela;
  \item suelos más realistas: otra función de correlación, rugosidad a varias escalas y
  distintos valores de rugosidad.
\end{itemize}
El nulo es estrecho (4 a 8° de $\Delta\phi$), así que también mediría cuánto se rellena y
con qué precisión hay que mantener $\Delta\phi$ en vuelo.

% ════════════════════════════════════════════════════════════════════
\section{Forma de combinar los canales}

\encorto{la suma con pesos iguales probablemente no es la mejor combinación.}
{de acuerdo; el Exp.~C2 lo confirma.}

Cuando la imagen bi tiene mucho menos clutter que la mono, sumarlas con el mismo peso
devuelve el clutter de la mono a la fusión. A 140°, la imagen bi sola supera a la fusión en
más de 10~dB de SCR. La suma de potencias (incoherente) ya la probamos: no necesita
calibrar la fase, pero rinde 2~dB menos de ISLR y 4.4~dB menos de contraste que la suma
coherente.

\paso{Qué propongo}
Comparar cuatro métodos en las mismas escenas (punto, objetos reales y clutter):
\begin{enumerate}
  \item suma coherente con pesos iguales (referencia);
  \item suma coherente con pesos por canal, elegidos para el mejor SCR o el menor ISLR;
  \item reconstrucción conjunta de los dos canales, que además puede estimar la fase entre
  ellos;
  \item suma de potencias (incoherente).
\end{enumerate}

% ════════════════════════════════════════════════════════════════════
\section{Calibración entre canales}

\encorto{suponer una única fase constante entre canales es el caso ideal; en la práctica
puede variar con el tiempo y la frecuencia.}
{nuestra hipótesis es que el desfase se mantiene constante durante la colecta. Podemos
simular el escenario variable, pero después de una evaluación previa que diga si hace falta.}

La fusión coherente necesita conocer el desfase entre el canal mono y el bi. Nuestra
intuición es que ese desfase es el mismo a lo largo de toda la colecta biestática: los dos
canales reciben el mismo pulso y el hardware no cambia durante el vuelo. Si es así, basta
con estimarlo una vez (por ejemplo, sobre el lóbulo principal de un blanco conocido), como
ya hace el relatorio.

El escenario que usted plantea es más general:
\[
  \varphi(f,t) = \varphi_0 + 2\pi f\,\Delta\tau(t) + H(f).
\]
Aquí $\varphi_0$ es el desfase constante (nuestra hipótesis). $\Delta\tau(t)$ es un
retardo que cambia durante la colecta, por deriva de reloj o error de posición del Rx.
$H(f)$ es la diferencia de respuesta en frecuencia entre las dos cadenas de recepción.

\paso{Qué propongo: primero una evaluación previa}
Antes de simular el modelo completo, estimar cuánto puede valer cada término con el
hardware que usaremos y compararlo con la tolerancia que ya conocemos. Las simulaciones del
relatorio dicen que la fusión aguanta un error de fase residual de hasta unos 30°, lo que a
425~MHz equivale a unos 0.2~ns de retardo o unos 6~cm de error en la longitud del trayecto.
\begin{enumerate}
  \item \textbf{Deriva de reloj:} con la estabilidad de los osciladores y la duración de la
  colecta, calcular el retardo acumulado.
  \item \textbf{Posición del Rx:} con la precisión del RTK, calcular el error de trayecto.
  \item \textbf{Cadenas de recepción:} medir o estimar la diferencia de respuesta en la
  banda.
\end{enumerate}

Si los tres quedan por debajo de la tolerancia, la hipótesis de fase constante es
suficiente y lo dejamos documentado. Si alguno la supera, simulamos sólo ese término y
diseñamos su corrección, por ejemplo estimando el retardo pulso a pulso con la señal
directa Tx$\to$Rx. Con datos reales, la fase de esa señal directa a lo largo del vuelo
permite comprobar la hipótesis de forma directa.

% ════════════════════════════════════════════════════════════════════
\section{Aperturas parciales y suelos húmedos}

\encorto{los resultados usan la hélice completa y no tratan bien los suelos húmedos.}
{las aperturas parciales faltan. La pérdida en suelo húmedo es un límite del sistema, no
de la trayectoria.}

Todos los resultados usan las dos vueltas completas de la hélice. Con un arco parcial la
imagen cambia de forma, y la mejor separación $\Delta\phi$ puede ser otra.

En suelo húmedo la pérdida de ida y vuelta hasta un blanco a 2~m (4~m dentro del suelo) es
muy grande y afecta por igual a la mono y al bi:

\begin{center}
\begin{tabular}{@{}lcc@{}}
\toprule
Suelo & Atenuación [dB/m] & Pérdida ida y vuelta a 2~m [dB] \\
\midrule
S1 & 7.7 & $\approx31$ \\
S2 & 13.6 & $\approx54$ \\
S3 & 27.7 & $\approx111$ \\
\bottomrule
\end{tabular}
\end{center}

\paso{Qué propongo}
\begin{itemize}
  \item Repetir los Exp.~B y F con arcos de 45°, 90°, 180° y 360°.
  \item Calcular, para cada suelo, la profundidad máxima a la que el blanco todavía
  supera al clutter (SCR~$>0$), y cómo crece al bajar la frecuencia.
\end{itemize}
Así acotamos dónde vale la propuesta, en lugar de presentarla como general.

% ════════════════════════════════════════════════════════════════════
\section{Prioridad propuesta}

\begin{enumerate}
  \item Puntos 1 y 3, porque definen la trayectoria y la forma de combinar los canales.
  \item Punto 2, que es lo que justificaría la novedad del artículo.
  \item Puntos 4 y 5, más acotados, en paralelo. El punto 4 empieza por la evaluación
  previa.
\end{enumerate}

% ════════════════════════════════════════════════════════════════════
\section{Preguntas para la próxima reunión}

\begin{enumerate}
  \item ¿Qué podemos publicar, y en qué congreso o revista?
  \item ¿Cuáles son los pasos para finalizar la investigación, además de la colecta y el
  procesamiento de datos reales?
  \item ¿Podemos definir el objetivo, la contribución y los resultados del cuerpo de la
  tesis?
\end{enumerate}

\vspace{0.5cm}
Saludos,\\
David

\end{document}
